\subsection{Extended compiler ledger}

\begin{table}[H]
\centering
\scriptsize
\sisetup{round-mode=none, group-minimum-digits=20}
\setlength{\tabcolsep}{2pt}
\renewcommand{\arraystretch}{0.96}
\begin{tabularx}{\linewidth}{@{}>{\raggedright\arraybackslash}p{0.15\linewidth}>{\raggedright\arraybackslash}p{0.33\linewidth}>{\raggedleft\arraybackslash}p{0.15\linewidth}Y@{}}
\toprule
Row & TFPT source or target & TFPT value & Status note \\
\midrule
$m_\gamma$ & exact unbroken $U(1)_{\mathrm{em}}$ direction & $0$ & exact structural row \\
$m_g$ & exact color gauge direction before confinement & $0$ & exact structural row \\
$W^\pm$ & boundary electroweak anchor at $M_Z$ & \makecell[r]{\num{80.369}\\GeV} & anchor-dependent comparison row \\
$Z$ & boundary electroweak anchor at $M_Z$ & \makecell[r]{\num{91.188}\\GeV} & anchor-dependent comparison row \\
$H$ & boundary Higgs anchor & \makecell[r]{\num{125.25}\\GeV} & anchor-dependent comparison row \\
\makecell[l]{$\hat s_Z^2$\\$(M_Z)$} & \makecell[l]{local gauge closure plus\\seam-even matching} & \num{0.23116113} & matching-branch comparison row \\
\makecell[l]{$\alpha_s$\\$(M_Z)$} & \makecell[l]{local gauge closure plus\\singlet Wilson matching} & \num{0.11796213} & matching-branch comparison row \\
$\Sigma m_\nu$ & oscillation-inversion normal-ordering branch sum & \makecell[r]{\num{5.8764e-2}\\eV} & oscillation-derived branch row \\
$M_R$ & seesaw-side scale cross-check / stage-atlas node & \makecell[r]{\num{1.31e15}\\GeV} & cross-check rather than hard derivation \\
$\Sigma m_\nu^{\mathrm{ss}}$ & geometric seesaw cross-check from $\vgeo$ and $M_R$ & \makecell[r]{\num{5.764e-2}\\eV} & separate from the operational branch row \\
\bottomrule
\end{tabularx}
\caption{Supplementary gauge, Higgs, and neutrino rows for the extended compiler ledger. Exact zero rows, matching rows, anchor rows, and branch / cross-check rows are intentionally separated.}
\label{tab:extended-compiler-ledger}
\end{table}

\begin{table}[H]
\centering
\scriptsize
\begin{tabularx}{\linewidth}{@{}>{\raggedright\arraybackslash}p{0.11\linewidth}>{\centering\arraybackslash}p{0.16\linewidth}>{\raggedleft\arraybackslash}p{0.16\linewidth}Y@{}}
\toprule
Field & Word-length data $L_{f,j}$ & Residual factor $\Lambda_{f,j}$ & Status note \\
\midrule
$e$ & $8$ & \num{0.466} & bounded $O(1)$ transport residual \\
$\mu$ & $5$ & \num{1.085} & bounded $O(1)$ transport residual \\
$\tau$ & $3$ & \num{0.899} & bounded $O(1)$ transport residual \\
$u$ & $7$ & \num{0.440} & bounded $O(1)$ transport residual \\
$d$ & $7$ & \num{0.936} & bounded $O(1)$ transport residual \\
$s$ & $5$ & \num{0.943} & bounded $O(1)$ transport residual \\
$c$ & $3$ & \num{0.646} & bounded $O(1)$ transport residual \\
$b$ & $2$ & \num{0.480} & bounded $O(1)$ transport residual \\
$t$ & $0$ & \num{0.986} & bounded $O(1)$ transport residual \\
\bottomrule
\end{tabularx}
\caption{Diagonal transport residuals in the Yukawa-factorization reading $y_{f,j}\approx \lambda_Y^{L_{f,j}}\Lambda_{f,j}$. The residual factors stay order one, which supports the transport grammar but does not remove the remaining theorem gap.}
\label{tab:yukawa-residuals}
\end{table}

\subsection{Novelty boundary and literature positioning}

The paper uses standard tools from relative spectral geometry, APS index theory, transport/resolvent language, and conventional RG comparison. The claim of novelty is not that these background ingredients are new in isolation, but that the minimal $3+2$ carrier packet, the compression identities, and the carrier-first compiler architecture are organized into one tightly overdetermined program. Conversely, the manuscript does not claim novelty for standard benchmark conventions, nor does it claim that appendix-level information-theoretic or quasicrystal readings are part of the hard core.

\begin{thebibliography}{99}

\bibitem{aps1975}
M.~F.~Atiyah, V.~K.~Patodi, and I.~M.~Singer,
\newblock \emph{Spectral asymmetry and Riemannian geometry. I},
\newblock Math. Proc. Cambridge Philos. Soc. \textbf{77} (1975), 43--69.

\bibitem{chamseddineConnes1997}
A.~H.~Chamseddine and A.~Connes,
\newblock \emph{The spectral action principle},
\newblock Commun. Math. Phys. \textbf{186} (1997), 731--750.

\bibitem{codata2022}
P.~J.~Mohr, D.~B.~Newell, and B.~N.~Taylor,
\newblock \emph{CODATA recommended values of the fundamental physical constants: 2022 update},
\newblock NIST / CODATA reference set, \url{https://physics.nist.gov/cuu/Constants/}, accessed March 2026.

\bibitem{pdg2025}
S.~Navas et al.\ (Particle Data Group),
\newblock \emph{Review of Particle Physics},
\newblock Phys.\ Rev.\ D \textbf{110} (2024), 030001; 2025 online update, \url{https://pdg.lbl.gov/2025/}, accessed March 2026.

\bibitem{nufit2025}
NuFIT Collaboration,
\newblock \emph{NuFIT global analysis of neutrino oscillation data},
\newblock website snapshot based on data available through November 2025, \url{https://www.nu-fit.org/}, accessed March 2026.

\bibitem{planck2018}
Planck Collaboration,
\newblock \emph{Planck 2018 results. VI. Cosmological parameters},
\newblock Astron. Astrophys. \textbf{641} (2020), A6.

\bibitem{minamiKomatsu2020}
Y.~Minami and E.~Komatsu,
\newblock \emph{New extraction of the cosmic birefringence from the Planck 2018 polarization data},
\newblock Phys. Rev. Lett. \textbf{125} (2020), 221301.

\bibitem{eskilt2023}
J.~R.~Eskilt et al.,
\newblock \emph{Cosmoglobe DR1 results. II. Constraints on isotropic cosmic birefringence from reprocessed WMAP and Planck LFI data},
\newblock Astron. Astrophys. \textbf{679} (2023), A144.

\bibitem{sullivan2025}
R.~M.~Sullivan, A.~Abghari, P.~Diego-Palazuelos, L.~T.~Hergt, and D.~Scott,
\newblock \emph{Planck PR4 (NPIPE) map-space cosmic birefringence},
\newblock arXiv:2502.07654.

\end{thebibliography}

\end{document}
